%\section{Concluding remarks}
\label{section: concluding remarks}
\section*{Conclusion}
This thesis studies stochastic impulse control problems on a finite time horizon, with applications to foreign exchange intervention and dynamic portfolio allocation. We perform analysis of the problem from two complementary perspectives. First, a probabilistic approach based on the dynamic programming principle is used to recursively solve the value function with iterated optimal stopping times. Second, a partial integro-differential equation approach characterises the value function as the viscosity solution of the quasi-variational inequality. 

The numerical part of the thesis connects reinforcement learning theory with the impulse control problem. We show by means of numerical experiments that deep reinforcement learning can approximate impulse-control problems and compare the results against a finite-difference scheme in the foreign exchange rate example. In the portfolio allocation problem, risky asset parameters were estimated using market data and correlated equity price paths simulated using multivariate subordination. The Proximal Policy Optimisation algorithm was then combined with Self-imitation Learning to trade and maximise portfolio value over a finite time horizon. The results indicate that the reinforcement learning algorithm effectively learns a trading strategy and achieves superior performance compared to a naive buy‑and‑hold portfolio.

\section*{Future research}
A first direction for future research is to extend the impulse-control problem beyond risk-neutral objectives. For instance, in the dynamic portfolio allocation experiment, the investor maximises expected discounted terminal wealth, while risk is controlled indirectly through the choice of equity universe, position limits, no-shorting, no-borrowing, and solvency constraints. While such constraints may resemble risk limits used by financial intermediaries\footnote{I.e. banks, insurance companies, and asset managers.}, they do not directly include risk preferences in the objective. A more realistic formulation could instead use a convex risk measure or a risk-sensitive criterion. For instance, \cite[Chapter 4.5.2]{oksendal2007applied} discusses the convex dual representation of convex risk measures, which provides a possible starting point for robust impulse-control formulations. Such dual representations can also be related to zero-sum impulse games, where the investor and market can be interpreted as competing decision-makers. Theory on zero-sum games is treated in \cite{azimzadeh2019zero, cosso2013stochastic}. Furthermore, finite-horizon risk-sensitive criteria for jump models are studied in \cite{davis2009jump, davis2011jump, davis2013jump, das2018risk}, while neural-network approaches can be found in \cite{bielecki2022risk, buehler2019deep, buehler2022deep, enders2024risk, fei2021exponential, noorani2025risk}.

A second direction concerns numerical schemes for HJBQVIs with jump-diffusion dynamics. The explicit-impulse scheme used in this thesis provides a useful benchmark in the one-dimensional diffusion setting, but it does not solve the QVI with generators containing nonlocal terms. An extension would be to split the PIDE into local and nonlocal parts and use an implicit-explicit discretisation of the differential operator. However, because the resulting QVI contains a nonlocal intervention operator, the scheme would require its own convergence proof. In particular, one would need to establish consistence, stability, monotonicity, and convergence for the QVI scheme. Starting points are the Barles-Souganidis convergence result \cite{barlessouganidis1991convergence}, the numerical impulse-control schemes in \cite{azimzadeh2016weakly, azimzadeh2018impulsecontrolfinancenumerical}, and finite-difference methods for PIDEs found in \cite{biswas2010difference}.

A third direction is to study reinforcement learning algorithms in a more systematic fashion. The DDQN algorithm is well-suited for problems with discrete impulse-sets, but it is sensitive to hyperparameters. Future work should therefore include a systematic hyperparameter study. In addition, other DRL algorithms could be considered, such as the Soft actor-critic (SAC), the Advantage Actor-Critic (A2C), and the Asynchronous Advantage Actor-Critic (A3C) algorithms \cite{A3C16, haarnoja2018soft}, while deep stopping time approaches also look promising, see \cite{becker2019deep, han2023new, lauriere2025deep}.

Finally, further work is needed on the portfolio experiment and the economic scenario generator. Currently, we rely on a correlated Variance‑Gamma model, calibrated partly to option prices and partly to historical stock data. While this approach captures jump behaviour and non‑Gaussian return features of equities, the resulting trading policy may still be sensitive to the estimated drift parameters. Future work should therefore include backtesting and stress testing across different market regimes, together with more robust benchmark portfolios. For example, stress tests based on multidimensional t-copulas, Archimedean copulas, or even Lévy copulas could provide a strong assessment of robustness based on tail dependence of stock prices, co‑jumps, and market‑wide shocks \cite{mcneil2015quantitative, tankov2004simulation, KALLSEN20061551, GROTHE2015182}. From a stock price modelling perspective, regime‑switching models \cite{korn2017stochastic, das2018risk, wu2025recursive} and models with mean‑reverting drift may better capture long‑term market dynamics than the currently used constant‑drift coefficients. In addition, enriching the state representation with features such as market sentiment \cite{MAO2024102048}, regime indicators, or volatility measures may lead to more robust trading policies.